\documentclass[10pt]{article}
\usepackage[T1]{fontenc}
\usepackage[utf8]{inputenc}
\usepackage[french]{babel}
\usepackage{amsmath,amssymb,mathrsfs,mathpazo}
\usepackage{geometry}
\geometry{a4paper,margin=2cm}
\usepackage{graphicx,array,tabularx,tabularray,colortbl,multirow,multicol,booktabs,enumitem,pifont,ulem,fancybox,diagbox,cancel}
\usepackage{pstricks,pstricks-add,pst-plot,pst-tree,pst-node,pst-text,pst-func,pst-math,pst-bspline,pst-3dplot,pst-eucl,pst-solides3d}
\let\pscalculate\undefined
\usepackage{pst-bezier}
\usepackage{tikz,pgfplots,tkz-tab}
\usetikzlibrary{arrows,arrows.meta,patterns,calc,positioning,trees,decorations.markings,shapes.geometric,shapes.misc,shapes.arrows,intersections,decorations.pathmorphing,decorations.pathreplacing,backgrounds,shadows,babel}
\pgfplotsset{compat=1.18}
\usepgfplotslibrary{fillbetween}
\usepackage[np]{numprint}
\usepackage{listings,makecell,pas-tableur,twoopt}
\let\barre\undefined
\usepackage[pstricks]{tablvar}
\definecolor{darkgreen}{RGB}{0,100,0}
\providecommand{\red}{\color{red}}
\providecommand{\blue}{\color{blue}}
\providecommand{\green}{\color{green}}
\providecommand{\black}{\color{black}}
\providecommand{\white}{\color{white}}
\providecommand{\gray}{\color{gray}}
\providecommand{\cyan}{\color{cyan}}
\providecommand{\magenta}{\color{magenta}}
\providecommand{\yellow}{\color{yellow}}
\providecommand{\violet}{\color{violet}}
\providecommand{\orange}{\color{orange}}
\providecommand{\euro}{\texteuro}
\providecommand{\EUR}{\texteuro}
\providecommand{\ieme}{\textsuperscript{e}}
\providecommand{\degres}{\ensuremath{^\circ}}
\providecommand{\degree}{\ensuremath{^\circ}}
\providecommand{\vv}[1]{\overrightarrow{#1}}
\providecommand{\Coord}[1]{#1}
\providecommand{\decosix}{}
\providecommand{\starredbullet}{\ensuremath{\bullet}}
\pagestyle{empty}
\setlength{\parindent}{0pt}
\usepackage{adjustbox}\DeclareUnicodeCharacter{2236}{:}\DeclareUnicodeCharacter{2190}{\ensuremath{\leftarrow}}
\providecommand{\R}{\mathbb{R}}
\providecommand{\N}{\mathbb{N}}
\providecommand{\Z}{\mathbb{Z}}
\providecommand{\Q}{\mathbb{Q}}
\providecommand{\D}{\mathbb{D}}
\providecommand{\C}{\mathbb{C}}
\providecommand{\e}{{\mathrm{e}}}
\providecommand{\vect}[1]{\overrightarrow{#1}}
\providecommand{\vectt}[1]{\overrightarrow{\mathrm{#1}}}
\providecommand{\vvt}[1]{\overrightarrow{\mathrm{#1}}}
\providecommand{\np}[2][]{{#2~#1}}
\providecommand{\vv}[1]{{\overrightarrow{#1}}}
\renewcommand{\d}{{\mathrm{d}}}
\providecommand{\degres}{{^\circ}}
\providecommand{\degre}{{^\circ}}
\providecommand{\blue}{\color{blue}}
\providecommand{\red}{\color{red}}
\providecommand{\enskip}{\quad}
\providecommand{\llbracket}{\lbrack\!\lbrack}
\providecommand{\rrbracket}{\rbrack\!\rbrack}
\providecommand{\euro}{{\text{€}}}
\providecommand{\no}{{\text{n°}}}
\providecommand{\interval}[2]{{[#1;#2]}}
\providecommand{\intervFO}[2]{{[#1;#2[}}
\providecommand{\dotfill}{\ldots}

\providecommand{\ds}{\displaystyle}
\providecommand{\textding}[1]{#1}
\providecommand{\Oij}{\(\mathrm{O};\vec{\imath},\vec{\jmath}\)}
\providecommand{\Oijk}{\(\mathrm{O};\vec{\imath},\vec{\jmath},\vec{k}\)}
\providecommand{\Ouv}{\(\mathrm{O};\vec{u},\vec{v}\)}
\providecommand{\intervFF}[2]{\ensuremath{[#1;#2]}}
\renewcommand{\euro}{\ifmmode\text{\texteuro}\else\texteuro\fi}
\DeclareUnicodeCharacter{00B0}{\ifmmode^{\circ}\else\textdegree\fi}
\DeclareUnicodeCharacter{20AC}{\ifmmode\text{\texteuro}\else\texteuro\fi}
% Énoncés : Éducation nationale. Transcription et corrigés : APMEP, Denis Vergès et contributeurs. Découpage et classement : CoopMaths, Sébastien Lozano, Éric Elter et contributeurs. Adaptation de la présentation : Maths Explorer.
% CC BY-SA : https://coopmaths.fr/www/about/
% https://coopmaths.fr/alea/static/bac/2021/tex/bac_2021_05_ameriquenord_1_cor.tex
\begin{document}
\noindent\textit{APMEP / CoopMaths. CC BY-SA. Présentation adaptée par Maths Explorer.}\par\medskip
\textbf{Commun à tous les candidats}










\bigskip

\textbf{Partie A}

\medskip




















\begin{enumerate}
\item ~
\begin{center}
\pstree[treemode=R,nodesepB=3pt,levelsep=2.75cm]{\TR{}}
{\pstree{\TR{$D$~~}\taput{0,08}}
	{\TR{$T$}\taput{0,98}
	\Tr{$\overline{T}$}\tbput{0,02}
	}
\pstree{\TR{$\overline{D}$~~}\tbput{0,92}}
	{\TR{$T$}\taput{0,005}
	\Tr{$\overline{T}$}\tbput{0,995}
	}	
}
\end{center}

\item 
D'après la loi des probabilités totales :

$P(T) = P(D \cap T) + P\left(\overline{D} \cap T \right)$ : or

$P(D \cap T) = P(D) \times P_D(T) =  0,08 \times 0,98 = \np{0,0784}$ et 

$P\left(\overline{D} \cap T \right) = P\left(\overline{D}  \right) \times P_{\overline{D}}(T) = 0,92 \times 0,005 = \np{0,00460}$. Donc :

$P(T)  = \np{0,0784} + \np{0,0046} = 0,083$.
\item 
	\begin{enumerate}
		\item 
La probabilité conditionnelle $P_T(D) = \dfrac{P(T \cap D)}{P(T)} = \dfrac{P(D \cap T)}{P(T)} = \dfrac{\np{0,0784}}{0,083} \approx \np{0,9445}$, soit 0,945 au millième près.
		\item 
D'après la question précédente $0,945 < 0,95$, donc le test ne sera pas commercialisé.


	\end{enumerate}
\end{enumerate}

\bigskip

\textbf{Partie B}

\medskip



\medskip

\begin{enumerate}
\item 


	\begin{enumerate}
		\item 
		Quel que soit l'athlète choisi la probabilité que cet  athlète présente un test positif est $0,103$. La variable aléatoire $X$ suit donc une loi binomiale de paramètres $n = 5$ et $p = 0,103$.
		\item 
On sait que $E = n \times p = 5 \times 0,103 = 0,515$ : ceci montre que sur un grand nombre de contrôles, il y aura à peu près 1 athlète sur 10 contrôlé positif.
		\item 
La probabilité qu'aucun athlète ne soit contrôlé positif est :

$0,103^0 \times (1 - 0,103)^5 = 0,897^5 \approx \np{0,5807}$ soit environ 0,581 au millième près.

Donc la probabilité qu'au moins un des 5 athlètes contrôlés présente un test positif est :

$1 - 0,581$, soit 0,419 au millième près.
	\end{enumerate}
\item 
On a pour $n$ athlètes contrôlés, $P(X = 0) = 0,103^0 \times 0,897^n = 0,897^n$.

Il faut donc trouver $n$ tel que :

$1 - 0,897^n \geqslant 0,75 \iff 1 - 0,75 \geqslant 0,897^n \iff 0,25 \geqslant 0,897^n$.

La calculatrice donne le plus petit $n \in \N$ vérifiant cette inéquation : pour $n = 13$, on a $0,897^{13} \approx 0,243$.

Il faut donc contrôler 13 athlètes en moyenne pour en trouver un positif.
\end{enumerate}
\vspace{0,5cm}
\bigskip

\end{document}
